\chapter{Debugging}\label{ch:debugging}

\index{debugging}\index{debugger}\index{clat run --debug}

Tests tell you \emph{that} something broke.
The debugger tells you \emph{why}.
Lattice includes a built-in interactive debugger that lets you pause execution, inspect variables, step through code line by line, and evaluate expressions---all without leaving the terminal.
For developers who prefer a graphical environment, the debugger also speaks the Debug Adapter Protocol (DAP), enabling first-class integration with VS Code and other editors.
And for the trickiest bugs---the ones where a value changed ``somewhere'' but you don't know where---Lattice's value history functions let you track mutations over time.

%% ─────────────────────────────────────────────────────────────
\section{The \texttt{--debug} Flag and the CLI Debugger}\label{sec:debug-flag}
\index{--debug flag}\index{CLI debugger}

To start a program under the debugger, add the \lstinline|--debug| flag:

\begin{lstlisting}[language=Lattice, caption={Launching the debugger}]
// $ clat run --debug server.lat
// Lattice debugger attached.
// Type 'help' for commands.
//
// Stopped at line 1 in <script>
//    1 | let host = "localhost"
//
// (dbg)
\end{lstlisting}

The program pauses at the very first line.
You're now in the \emph{debug REPL}---an interactive prompt where you can inspect state, set breakpoints, and control execution.

\begin{defnbox}{The Debug REPL}
The debug REPL is a command loop that activates whenever execution is paused---either at startup in step mode, at a breakpoint, or after a step command completes.
The prompt \lstinline|(dbg)| indicates that the program is paused and waiting for your instruction.
All debugger commands are processed by the \lstinline|debugger_check()| function in \filename{src/debugger.c}.
\end{defnbox}

\subsection{The Help Command}\label{sec:debug-help}

Type \lstinline|help| (or \lstinline|h|) at the debug prompt to see all available commands:

\begin{lstlisting}[language=Lattice, caption={Debugger help output}]
// (dbg) help
// Debugger commands:
//   s, step           Step into (execute one line, enter calls)
//   n, next           Step over (execute one line, skip calls)
//   o, out, finish    Step out (run until current function returns)
//   c, continue       Continue until next breakpoint
//   l, locals         Show local variables
//   bt, stack         Show call stack / backtrace
//   list              Show source context around current line
//   b <line|func>     Set breakpoint at line or function
//   b <N> if <expr>   Set conditional breakpoint
//   d <id>            Delete breakpoint by ID
//   p <var>           Print variable value
//   eval <expr>       Evaluate expression
//   watch <expr>      Add watch expression
//   unwatch <id>      Remove watch expression
//   info b            List breakpoints
//   info w            List watch expressions
//   q, quit           Exit debugger
//   h, help           Show this help
\end{lstlisting}

\subsection{Initial Breakpoint with \texttt{--break}}\label{sec:break-line-flag}
\index{--break flag}

If you know roughly where the bug is, you can jump straight to a specific line instead of stepping from the beginning:

\begin{lstlisting}[language=Lattice, caption={Starting with a breakpoint}]
// $ clat run --break 42 server.lat
// Lattice debugger attached.
// Breakpoints: line 42
//
// Breakpoint 1 at line 42
//    42 | let response = handle_request(request)
//
// (dbg)
\end{lstlisting}

The \lstinline|--break| flag implies \lstinline|--debug|.
The program runs at full speed until it hits line 42, then pauses.

\subsection{Viewing Source Context}\label{sec:debug-list}
\index{debugger!list command}

The \lstinline|list| command shows the source code surrounding the current line, with an arrow pointing to where execution is paused:

\begin{lstlisting}[language=Lattice, caption={Viewing source context}]
// (dbg) list
//    37 |     let method = request.get("method")
//    38 |     let path = request.get("path")
//    39 |
//    40 |     if method == "GET" {
//    41 |         return handle_get(path)
// -->42 |     let response = handle_request(request)
//    43 |     } else if method == "POST" {
//    44 |         return handle_post(path, body)
//    45 |     }
//    46 |
//    47 |     return error("unsupported method: " + method)
\end{lstlisting}

The debugger loads the source file at startup (via \lstinline|debugger_load_source()| in \filename{src/debugger.c}) so it can display context around the current execution point.

\subsection{Inspecting Variables}\label{sec:debug-inspect}
\index{debugger!variable inspection}

The \lstinline|p| (or \lstinline|print|) command shows the value of a variable:

\begin{lstlisting}[language=Lattice, caption={Printing a variable}]
// (dbg) p request
// request = {method: "GET", path: "/api/users", headers: {...}}
//
// (dbg) p method
// method = "GET"
\end{lstlisting}

The \lstinline|locals| command (or \lstinline|l|) shows \emph{all} local variables in the current scope:

\begin{lstlisting}[language=Lattice, caption={Viewing all local variables}]
// (dbg) locals
//   request = {method: "GET", path: "/api/users", ...}
//   method = "GET"
//   path = "/api/users"
\end{lstlisting}

Under the hood, the debugger reads local variable names from the chunk's \lstinline|local_names| array and their values from the VM's stack slots.
If a variable has no debug name (e.g., a compiler-generated temporary), it won't appear in the listing.

\subsection{Evaluating Expressions}\label{sec:debug-eval}
\index{debugger!expression evaluation}

The \lstinline|eval| (or \lstinline|e|) command evaluates an arbitrary Lattice expression in the context of the paused program:

\begin{lstlisting}[language=Lattice, caption={Evaluating expressions in the debugger}]
// (dbg) eval len(path)
// => 11
//
// (dbg) eval method == "GET"
// => true
//
// (dbg) eval path.split("/")
// => ["", "api", "users"]
\end{lstlisting}

This is one of the most powerful features of the debugger.
You can call functions, access fields, perform arithmetic---anything you could write in a normal Lattice expression.
The evaluator compiles the expression on the fly and executes it on the live VM (see \lstinline|debugger_eval_expr()| in \filename{src/debugger.c}).

\begin{warnbox}{Side Effects in Eval}
Expressions evaluated with \lstinline|eval| run on the real VM state.
If you evaluate something with side effects (e.g., \lstinline|eval items.push(99)|), those effects will persist when you continue execution.
Use \lstinline|eval| for inspection, not mutation---unless you know what you're doing.
\end{warnbox}

\subsection{The Call Stack}\label{sec:debug-backtrace}
\index{debugger!backtrace}\index{debugger!call stack}

The \lstinline|bt| (or \lstinline|stack| or \lstinline|backtrace|) command shows the complete call chain:

\begin{lstlisting}[language=Lattice, caption={Viewing the call stack}]
// (dbg) bt
// Backtrace:
//   #0 [line 12] in <script>
//   #1 [line 25] in process_request()
//   #2 [line 42] in handle_request()  <-- current
\end{lstlisting}

Frame \lstinline|#0| is the bottom of the stack (the top-level script), and the frame marked \lstinline|<-- current| is where execution is paused.
This helps you understand how you got to the current line---especially useful when the same function is called from multiple places.

%% ─────────────────────────────────────────────────────────────
\section{Breakpoints: Line, Function, Conditional}\label{sec:breakpoints}
\index{breakpoints}

Breakpoints tell the debugger ``stop here.''
Lattice supports three kinds: line breakpoints, function breakpoints, and conditional breakpoints.

\subsection{Line Breakpoints}\label{sec:bp-line}
\index{breakpoints!line}

Set a breakpoint at a specific line number:

\begin{lstlisting}[language=Lattice, caption={Setting a line breakpoint}]
// (dbg) b 15
// Breakpoint 1 set at line 15
//
// (dbg) c
// ... program runs ...
// Breakpoint 1 at line 15
//    15 | let total = calculate_total(items)
//
// (dbg)
\end{lstlisting}

The debugger assigns each breakpoint an auto-incrementing ID (starting at 1).
When execution hits the breakpoint, the debugger shows the breakpoint number and line.

\subsection{Function Breakpoints}\label{sec:bp-function}
\index{breakpoints!function}

If you don't know the line number but know the function name, set a function breakpoint:

\begin{lstlisting}[language=Lattice, caption={Setting a function breakpoint}]
// (dbg) b calculate_total
// Breakpoint 2 set on function calculate_total()
//
// (dbg) c
// ... program runs ...
// Breakpoint 2 at calculate_total(), line 28
//    28 | fn calculate_total(items: Array) -> Float {
//
// (dbg)
\end{lstlisting}

Function breakpoints trigger whenever the named function is entered, regardless of which line called it.
The debugger detects function entry by comparing the current frame count against the previous check---if the frame count increased and the new frame's function name matches, the breakpoint fires (see the function breakpoint check in \lstinline|debugger_check()|).

\subsection{Conditional Breakpoints}\label{sec:bp-conditional}
\index{breakpoints!conditional}

A conditional breakpoint only pauses when a Lattice expression evaluates to a truthy value.
The syntax is \lstinline|b <line|func> if <expression>|:

\begin{lstlisting}[language=Lattice, caption={Setting conditional breakpoints}]
// (dbg) b 15 if total > 1000.0
// Breakpoint 3 set at line 15 if total > 1000.0
//
// (dbg) b validate if input == ""
// Breakpoint 4 set on function validate() if input == ""
\end{lstlisting}

The condition is evaluated every time the breakpoint location is reached.
If the condition is falsy, execution continues without pausing.
If the condition throws an error, the debugger warns you and breaks anyway (so you can fix the condition).

\begin{tipbox}{Use Conditional Breakpoints for Loops}
If you're debugging a loop that runs thousands of times but only the 500th iteration is wrong, a conditional breakpoint is far more effective than stepping through all 500 iterations:
\begin{lstlisting}[language=Lattice]
// (dbg) b 20 if index == 500
\end{lstlisting}
\end{tipbox}

\subsection{Managing Breakpoints}\label{sec:bp-manage}
\index{breakpoints!management}

List all breakpoints with \lstinline|info b| (or \lstinline|info breakpoints|):

\begin{lstlisting}[language=Lattice, caption={Listing breakpoints}]
// (dbg) info b
// Num  Type      Enb  What
// 1    line      y    line 15  (hits: 3)
// 2    function  y    calculate_total()  (hits: 1)
// 3    line      y    line 15  if total > 1000.0  (hits: 0)
\end{lstlisting}

The table shows each breakpoint's ID, type (line or function), whether it's enabled (\lstinline|y|/\lstinline|n|), what it's set on, any condition, and how many times it has been hit.

Remove a breakpoint by its ID:

\begin{lstlisting}[language=Lattice, caption={Deleting a breakpoint}]
// (dbg) d 2
// Breakpoint 2 removed
\end{lstlisting}

%% ─────────────────────────────────────────────────────────────
\section{Step-Into, Step-Over, Step-Out}\label{sec:stepping}
\index{stepping}\index{step-into}\index{step-over}\index{step-out}

The three stepping commands are the bread and butter of interactive debugging.
They control how much code executes before the debugger pauses again.

\subsection{Step-Into (\texttt{s})}\label{sec:step-into}

The \lstinline|s| (or \lstinline|step|) command executes one line.
If that line contains a function call, the debugger enters the function and pauses at its first line:

\begin{lstlisting}[language=Lattice, caption={A simple program to step through}]
fn add_tax(price: Float, rate: Float) -> Float {
    let tax = price * rate
    return price + tax
}

let total = add_tax(100.0, 0.08)
print(total)
\end{lstlisting}

\begin{lstlisting}[language=Lattice, caption={Stepping into a function call}]
// Stopped at line 6 in <script>
//     6 | let total = add_tax(100.0, 0.08)
//
// (dbg) s
// Stopped at line 2 in add_tax()
//     2 |     let tax = price * rate
//
// (dbg) p price
// price = 100.0
//
// (dbg) s
// Stopped at line 3 in add_tax()
//     3 |     return price + tax
//
// (dbg) p tax
// tax = 8.0
\end{lstlisting}

Step-into follows the actual execution path.
If a line calls multiple functions, stepping enters the first one called.

\subsection{Step-Over (\texttt{n})}\label{sec:step-over}

The \lstinline|n| (or \lstinline|next|) command executes one line but treats function calls as single steps---it doesn't enter them:

\begin{lstlisting}[language=Lattice, caption={Stepping over a function call}]
// Stopped at line 6 in <script>
//     6 | let total = add_tax(100.0, 0.08)
//
// (dbg) n
// Stopped at line 7 in <script>
//     7 | print(total)
//
// (dbg) p total
// total = 108.0
\end{lstlisting}

The function \lstinline|add_tax| executed fully, but the debugger didn't pause inside it.
Step-over is the right choice when you trust the called function and want to see its result.

The implementation tracks the call depth: when \lstinline|next| is issued, the debugger records the current frame count and only pauses when the frame count is at or below that level \emph{and} the line has changed (see the \lstinline|next_mode| and \lstinline|next_depth| fields in the \lstinline|Debugger| struct).

\subsection{Step-Out (\texttt{o})}\label{sec:step-out}

The \lstinline|o| (or \lstinline|out| or \lstinline|finish|) command runs until the current function returns, then pauses in the caller:

\begin{lstlisting}[language=Lattice, caption={Stepping out of a function}]
// Stopped at line 2 in add_tax()
//     2 |     let tax = price * rate
//
// (dbg) o
// Stopped at line 7 in <script>
//     7 | print(total)
//
// (dbg) p total
// total = 108.0
\end{lstlisting}

Step-out is perfect when you've accidentally stepped into a function you don't care about.
Rather than stepping through every line, \lstinline|o| runs the rest of the function at full speed and returns you to the caller.

\subsection{Continue (\texttt{c})}\label{sec:continue}
\index{debugger!continue}

The \lstinline|c| (or \lstinline|continue|) command resumes execution at full speed until the next breakpoint is hit (or the program ends):

\begin{lstlisting}[language=Lattice, caption={Continuing execution}]
// (dbg) b 25
// Breakpoint 1 set at line 25
//
// (dbg) c
// ... program runs at full speed ...
// Breakpoint 1 at line 25
//    25 | let result = process(data)
//
// (dbg)
\end{lstlisting}

\begin{notebox}{Choosing the Right Step Command}
\begin{itemize}
  \item \textbf{Step-into (s)}: Use when you want to see what happens \emph{inside} a function call.
  \item \textbf{Step-over (n)}: Use when you trust the called function and want to move to the next line.
  \item \textbf{Step-out (o)}: Use when you're inside a function you don't care about and want to return to the caller.
  \item \textbf{Continue (c)}: Use when you've set breakpoints and want to skip ahead.
\end{itemize}
\end{notebox}

%% ─────────────────────────────────────────────────────────────
\section{Watch Expressions and \texttt{breakpoint()}}\label{sec:watches}
\index{watch expressions}\index{debugger!watch}

\subsection{Watch Expressions}\label{sec:watch-expr}

A watch expression is monitored automatically every time the debugger pauses.
Instead of typing \lstinline|p total| at every stop, add a watch:

\begin{lstlisting}[language=Lattice, caption={Adding watch expressions}]
// (dbg) watch total
// Watch 1: total
//
// (dbg) watch len(items)
// Watch 2: len(items)
\end{lstlisting}

Now every time the debugger pauses, it prints the current value of each watch expression:

\begin{lstlisting}[language=Lattice, caption={Watch expressions displayed at each stop}]
// (dbg) s
// Stopped at line 18 in process()
//    18 |     total = total + item.price
//
// Watch expressions:
//   [1] total = 342.50
//   [2] len(items) = 15
//
// (dbg)
\end{lstlisting}

Watch expressions can be any valid Lattice expression---not just variable names.
You can watch \lstinline|items[0].name|, \lstinline|total / len(items)|, or even a function call.

Manage watches with:

\begin{lstlisting}[language=Lattice, caption={Managing watch expressions}]
// (dbg) info w
//   [1] total = 342.50
//   [2] len(items) = 15
//
// (dbg) unwatch 2
// Watch 2 removed
\end{lstlisting}

\subsection{The \texttt{breakpoint()} Built-in}\label{sec:breakpoint-fn}
\index{breakpoint() function}

Sometimes the condition for pausing is too complex to express in a breakpoint condition string.
For these cases, Lattice provides the \lstinline|breakpoint()| built-in function.
Call it anywhere in your code, and the debugger pauses at that point if the program is running under \lstinline|--debug|:

\begin{lstlisting}[language=Lattice, caption={Using breakpoint() in code}]
fn process_order(order: Map) -> Map {
    let items = order.get("items")
    let total = 0.0

    for item in items {
        total = total + item.get("price")

        // Pause if we encounter a suspiciously high price
        if item.get("price") > 10000.0 {
            breakpoint()
        }
    }

    return order
}
\end{lstlisting}

When the debugger is not attached (normal \lstinline|clat run|), \lstinline|breakpoint()| is a no-op.
This means you can leave \lstinline|breakpoint()| calls in your code without affecting production behavior.

\begin{tipbox}{Programmatic Breakpoints}
\lstinline|breakpoint()| is especially useful for catching rare conditions that are hard to reproduce.
Rather than guessing which line to break on, embed the breakpoint directly in the condition logic:
\begin{lstlisting}[language=Lattice]
if is_error(result) {
    breakpoint()  // pause only when something went wrong
}
\end{lstlisting}
\end{tipbox}

%% ─────────────────────────────────────────────────────────────
\section{DAP and VS Code Integration}\label{sec:dap}
\index{DAP}\index{Debug Adapter Protocol}\index{VS Code}\index{debugger!VS Code}

The CLI debugger is powerful, but sometimes you want a graphical view: variable panes, inline source highlighting, and click-to-set breakpoints.
Lattice supports the Debug Adapter Protocol (DAP), the industry standard used by VS Code, Neovim (via nvim-dap), Emacs (via dap-mode), and other editors.

\subsection{What Is DAP?}\label{sec:what-is-dap}

The Debug Adapter Protocol, developed by Microsoft for VS Code, defines a JSON-based message protocol between an editor (the ``client'') and a debugger (the ``adapter'').
Instead of every editor implementing custom support for every language's debugger, DAP provides a common interface.
Lattice's DAP implementation lives in \filename{src/dap.c} and speaks DAP over stdin/stdout using Content-Length framed JSON messages.

\subsection{Starting DAP Mode}\label{sec:dap-start}

To launch the debugger in DAP mode:

\begin{lstlisting}[language=Lattice, caption={Launching DAP mode}]
// $ clat run --dap server.lat
\end{lstlisting}

In DAP mode, the debugger communicates via the Content-Length framed JSON protocol on stdin/stdout instead of the interactive REPL.
The handshake follows the DAP specification: the client sends \lstinline|initialize|, \lstinline|launch|, optionally \lstinline|setBreakpoints| and \lstinline|setFunctionBreakpoints|, and finally \lstinline|configurationDone|.

\subsection{VS Code Configuration}\label{sec:vscode-config}

To debug Lattice programs in VS Code, create a \filename{.vscode/launch.json} file in your project:

\begin{lstlisting}[language=Lattice, caption={VS Code launch.json for Lattice debugging}]
// {
//     "version": "0.2.0",
//     "configurations": [
//         {
//             "name": "Debug Lattice",
//             "type": "lattice",
//             "request": "launch",
//             "program": "server.lat",
//             "stopOnEntry": true
//         }
//     ]
// }
\end{lstlisting}

The \lstinline|"stopOnEntry": true| option tells the debugger to pause at the first line (equivalent to \lstinline|--debug| in CLI mode).
If you set it to \lstinline|false|, the program runs until it hits a breakpoint.

\subsection{DAP Capabilities}\label{sec:dap-capabilities}

Lattice's DAP adapter advertises the following capabilities during the initialize handshake (see \lstinline|dap_handshake()| in \filename{src/dap.c}):

\begin{longtable}{@{}lp{8cm}@{}}
\toprule
\textbf{Capability} & \textbf{Description} \\
\midrule
\endhead
Configuration Done & Supports the \lstinline|configurationDone| request \\
Function Breakpoints & Supports breakpoints on function names \\
Conditional Breakpoints & Supports breakpoint conditions \\
Evaluate for Hovers & Evaluates expressions when you hover in the editor \\
\bottomrule
\end{longtable}

\subsection{Features in the VS Code Debug View}\label{sec:vscode-features}

Once connected, you get the full VS Code debugging experience:

\begin{itemize}
  \item \textbf{Variables pane}: Shows all local variables and their values, organized by scope (Locals and Globals). Compound types (arrays, structs, maps) can be expanded to inspect their contents.
  \item \textbf{Call stack}: Displays the call chain from the top-level script down to the current function.
  \item \textbf{Breakpoint gutter}: Click in the gutter to set or remove line breakpoints. Right-click for conditional breakpoints.
  \item \textbf{Step controls}: The toolbar provides step-into, step-over, step-out, continue, and pause buttons.
  \item \textbf{Debug console}: Type expressions in the debug console to evaluate them, just like the CLI \lstinline|eval| command.
  \item \textbf{Output capture}: \lstinline|print()| output from your program appears in the Debug Console via DAP output events.
\end{itemize}

\begin{notebox}{Compound Variable Expansion}
When the Variables pane shows a struct, array, or map, VS Code can expand it to show individual fields or elements.
This works through DAP's variable reference system: each compound value gets a unique reference ID, and when you expand it, VS Code sends a \lstinline|variables| request with that reference.
The implementation in \filename{src/dap.c} uses a per-stop variable reference table (cleared on each stop event to avoid stale data).
\end{notebox}

\subsection{Using with Other Editors}\label{sec:dap-other-editors}

Any editor that supports DAP can connect to Lattice's debugger.
For Neovim with nvim-dap:

\begin{lstlisting}[language=Lattice, caption={Neovim nvim-dap configuration}]
// In your Neovim config (Lua):
// local dap = require("dap")
// dap.adapters.lattice = {
//     type = "executable",
//     command = "clat",
//     args = { "run", "--dap" }
// }
// dap.configurations.lattice = {{
//     name = "Debug Lattice",
//     type = "lattice",
//     request = "launch",
//     program = "${file}",
//     stopOnEntry = true,
// }}
\end{lstlisting}

The DAP protocol is identical regardless of the editor.
Any tool that implements the DAP client side will work with Lattice's adapter.

%% ─────────────────────────────────────────────────────────────
\section{Value History: \texttt{track()}, \texttt{history()}, \texttt{rewind()}}\label{sec:value-history}
\index{track()}\index{history()}\index{rewind()}\index{value history}\index{time travel debugging}

Some bugs aren't about the current state---they're about how the state got there.
A variable has the wrong value, but when did it change?
What was it before?
Lattice's value history functions let you answer these questions without setting breakpoints on every assignment.

\subsection{\texttt{track(variable)} --- Start Recording}\label{sec:track}
\index{track()}

The \lstinline|track()| function tells the runtime to record every value change for a named variable:

\begin{lstlisting}[language=Lattice, caption={Tracking a variable}]
flux counter = 0
track(counter)

counter = 10
counter = 20
counter = 30
\end{lstlisting}

After calling \lstinline|track(counter)|, every subsequent assignment to \lstinline|counter| is recorded as a snapshot.
The snapshot includes the value, its phase, the source line number, and the enclosing function name.

\begin{defnbox}{track(variable)}
\lstinline|track(variable)| enables value history tracking for the specified variable.
The argument is a variable name (not a string---the compiler automatically converts it to a string behind the scenes).
Tracking begins with an initial snapshot of the variable's current value.
Subsequent mutations are recorded as they happen.
\end{defnbox}

\subsection{\texttt{history(variable)} --- View the Timeline}\label{sec:history}
\index{history()}

Once you're tracking a variable, call \lstinline|history()| to retrieve its full timeline:

\begin{lstlisting}[language=Lattice, caption={Viewing variable history}]
flux temperature = 20.0
track(temperature)

temperature = 22.5
temperature = 25.0
temperature = 18.3

let timeline = history(temperature)
for entry in timeline {
    print(entry)
}
// Output:
// {phase: "flux", value: 20.0, line: 2, fn: nil}
// {phase: "flux", value: 22.5, line: 4, fn: nil}
// {phase: "flux", value: 25.0, line: 5, fn: nil}
// {phase: "flux", value: 18.3, line: 6, fn: nil}
\end{lstlisting}

Each entry in the timeline is a \lstinline|Map| with four keys:

\begin{longtable}{@{}lp{9cm}@{}}
\toprule
\textbf{Key} & \textbf{Description} \\
\midrule
\endhead
\lstinline|phase| & The phase of the variable at the time of recording (e.g., \lstinline|"flux"|, \lstinline|"fix"|, \lstinline|"let"|) \\
\lstinline|value| & A deep copy of the variable's value at that moment \\
\lstinline|line| & The source line number where the assignment occurred \\
\lstinline|fn| & The name of the enclosing function, or \lstinline|nil| for top-level code \\
\bottomrule
\end{longtable}

The \lstinline|phases()| function is a synonym for \lstinline|history()| that returns the same data. Use whichever name feels more natural.

\subsection{\texttt{rewind(variable, n)} --- Look Back in Time}\label{sec:rewind}
\index{rewind()}

The \lstinline|rewind()| function retrieves the value of a variable from \lstinline|n| steps ago:

\begin{lstlisting}[language=Lattice, caption={Rewinding variable state}]
flux score = 0
track(score)

score = 10
score = 25
score = 50
score = 100

print(rewind(score, 0))  // Output: 100 (current)
print(rewind(score, 1))  // Output: 50  (one step back)
print(rewind(score, 2))  // Output: 25  (two steps back)
print(rewind(score, 3))  // Output: 10  (three steps back)
print(rewind(score, 4))  // Output: 0   (initial value)
\end{lstlisting}

The parameter \lstinline|n| counts backward from the most recent snapshot.
\lstinline|rewind(score, 0)| returns the current value; \lstinline|rewind(score, 1)| returns the previous value, and so on.
If \lstinline|n| exceeds the number of recorded snapshots, \lstinline|rewind()| returns \lstinline|nil|.

The returned value is a \emph{deep copy}---modifying it won't affect the original variable or the history.

\subsection{Practical Example: Debugging a Running Total}\label{sec:history-example}

Let's use value history to find where a running total goes wrong:

\begin{lstlisting}[language=Lattice, caption={Debugging with value history}]
fn process_transactions(transactions: Array) -> Float {
    flux total = 0.0
    track(total)

    for tx in transactions {
        let amount = tx.get("amount")
        if tx.get("type") == "credit" {
            total = total + amount
        } else {
            total = total - amount
        }
    }

    // Something's wrong with the total.
    // Let's inspect the history:
    let timeline = history(total)
    for entry in timeline {
        print(format("Line {}: total = {} ({})",
            entry.get("line"),
            entry.get("value"),
            entry.get("phase")))
    }

    return total
}

let transactions = [
    {"type": "credit", "amount": 100.0},
    {"type": "credit", "amount": 50.0},
    {"type": "debit", "amount": 30.0},
    {"type": "credit", "amount": -20.0}
]

let result = process_transactions(transactions)
print("Final total: " + to_string(result))
// The history output shows exactly when the negative credit
// caused the unexpected subtraction.
\end{lstlisting}

By scanning the timeline, you can see exactly which transaction caused the total to deviate from expectations---and on which line.

\subsection{Combining Tracking with the Debugger}\label{sec:track-with-debug}

Value history and the interactive debugger complement each other.
You can set a breakpoint, then inspect the history at that point:

\begin{lstlisting}[language=Lattice, caption={Using track() with the CLI debugger}]
// $ clat run --debug analytics.lat
//
// (dbg) b 45
// Breakpoint 1 set at line 45
//
// (dbg) c
// Breakpoint 1 at line 45
//    45 |     return average
//
// (dbg) eval history(running_sum)
// => [{phase: "flux", value: 0, line: 12, fn: "compute_average"},
//     {phase: "flux", value: 42, line: 18, fn: "compute_average"},
//     {phase: "flux", value: 85, line: 18, fn: "compute_average"},
//     ...]
//
// (dbg) eval rewind(running_sum, 3)
// => 42
\end{lstlisting}

This is the closest thing to time-travel debugging: you can inspect what a variable's value was at any prior point in the program, without having to re-run it.

\begin{warnbox}{Memory Cost of Tracking}
Each tracked assignment creates a deep copy of the variable's value.
For large data structures in tight loops, this can consume significant memory.
Use \lstinline|track()| judiciously---enable it for the specific variables you need to debug, and remove it when you're done.
\end{warnbox}

%% ─────────────────────────────────────────────────────────────
\section*{Exercises}\label{sec:debugging-exercises}
\addcontentsline{toc}{section}{Exercises}

\begin{enumerate}
  \item Write a program that computes the factorial of a number using recursion. Run it under \lstinline|--debug| and step through the recursive calls using \lstinline|s|. Use \lstinline|bt| at the deepest recursion level to see the full call stack. Then use \lstinline|o| to step out one frame at a time.

  \item Create a program with a loop that processes an array of numbers. Set a conditional breakpoint that only triggers when the current number is negative. Use \lstinline|eval| to inspect the loop state when the breakpoint fires.

  \item Add \lstinline|track()| to a variable inside a sorting algorithm (e.g., insertion sort). After the sort completes, print the full \lstinline|history()| and use \lstinline|rewind()| to inspect the variable's value at three different points during execution.

  \item Set up a VS Code \filename{launch.json} for a Lattice project. Set breakpoints by clicking in the gutter, then step through the program using the toolbar controls. Expand a struct and an array in the Variables pane to see their fields.

  \item Write a program that uses \lstinline|breakpoint()| inside an error-handling path. Run it normally (without \lstinline|--debug|) to verify that \lstinline|breakpoint()| has no effect, then run it under the debugger to confirm it pauses at the right moment.

  \item Combine \lstinline|watch| expressions with \lstinline|track()| to debug a program that accumulates values in a loop. Set two watch expressions (one for the running total, one for the loop counter) and step through five iterations. Then use \lstinline|rewind()| to verify the total at each step.
\end{enumerate}

%% ─────────────────────────────────────────────────────────────
\section*{What's Next}\label{sec:debugging-next}
\addcontentsline{toc}{section}{What's Next}

You now have a complete toolkit for developing Lattice programs: formatting (\lstinline|clat fmt|), documentation (\lstinline|clat doc|), testing (\lstinline|clat test|), and debugging (\lstinline|--debug| and DAP).
In Part~X, we'll go under the hood and explore the internals of the Lattice runtime itself.
\Cref{ch:three-backends} starts with the three compilation backends---tree-walk interpreter, stack VM, and register VM---and explains how your Lattice code is transformed from source text into executable instructions.
